\section{Background: Resilience in Transportation Networks}
\label{sec:background}

\subsection{Two Dimensions of Resilience}

Transportation resilience research has historically addressed two distinct failure mechanisms. The first is structural: a network is resilient when it can reroute demand around a failed link without unacceptable cost. Berdica \citep{Berdica2002} defined link vulnerability as the susceptibility of a link to incidents that substantially reduce network serviceability; the key insight was that vulnerability is a network-level property, not a link-level one. A link is only vulnerable if its failure degrades the performance of the network as a whole. Mattsson and Jenelius \citep{Mattsson2015} extended this to a formal framework distinguishing link importance (the network degradation caused by link failure) from link exposure (the probability and intensity of failure-inducing events). A resilient network has either low importance for any individual link (high redundancy --- no single link matters much) or low exposure for important links.

The second dimension is environmental: a network is resilient when its physical infrastructure withstands the events that cause closure. Flood-resilient bridges, snowshed protection on mountain passes, hurricane-resistant barrier walls, and wildfire-hardened drainage systems all address environmental exposure by reducing the probability and duration of closure events, conditional on an event occurring.

These two dimensions are often treated as parallel investment priorities: federal resilience programs fund both alternate route development (structural) and infrastructure hardening (environmental). The ROUTE rubric operationalizes them as two separate dimensions: B1 measures structural isolation, and D1 measures environmental vulnerability.

\subsection{The B1 Dimension: Structural Isolation}

B1 is defined as the inverse of alternate-route adequacy. A corridor scores B1 = 0 when closure causes negligible network impact: adequate alternate routes exist at interstate standard, with detour distances below 100 additional miles and no freight-relevant grade restrictions. A corridor scores B1 = 10 when closure produces a national-scale freight disruption with no interstate-standard alternate: rerouting requires 200+ additional miles over routes not maintained for commercial vehicle traffic in relevant weather conditions.

The B1 score is operationalized using the ROUTE corpus methodology \citep{ROUTE_B1}: for each corridor segment, the nearest alternate route is identified from TIGER/Line geometry, its interstate-standard adequacy is assessed (load ratings, grade restrictions, winter maintenance), and the detour penalty (miles, hours, freight cost) is computed using ATRI truck-hour operating costs \citep{ATRI_costs2024}. The result is a per-corridor B1 score that reflects how much of a single-point-of-failure role the corridor plays in the national network.

Key B1 = 10 cases in the corpus are mountain passes and high-desert crossings where no parallel interstate exists: Donner Pass (I-80 Sierra Nevada), Snoqualmie Pass (I-90 Washington Cascades), and Homestake Pass (I-90 Montana Continental Divide). In each case, the alternate is a US highway not maintained for commercial freight during winter closure conditions --- the same conditions that trigger closure of the primary corridor.

\subsection{The D1 Dimension: Environmental Vulnerability}

D1 is a composite climate vulnerability score incorporating three exposure components: (1) FEMA National Flood Hazard Layer (NFHL) Special Flood Hazard Area miles as a share of total corridor miles, reflecting baseline flood exposure probability; (2) mean annual closure events from FHWA incident records and Caltrans/WSDOT/LADOTD operational logs \citep{Caltrans2023, FHWA_freight2023}; and (3) a climate trajectory adjustment reflecting whether the 2050 exposure under NOAA intermediate sea-level-rise and precipitation-intensity scenarios is substantially higher than the 2024 baseline \citep{FEMA_NFHL}.

A corridor scores D1 = 0 when it has no FHWA SFHA miles, near-zero historical closure frequency, and no projected increase in 2050 exposure. It scores D1 = 10 when a majority of the corridor lies in high-hazard flood zones, annual closure frequency is high (50+ events/year), and 2050 projections show substantially elevated exposure. The score is signed for vulnerability, not resilience: higher D1 is worse.

\subsection{Why Compound Exposure Is Not Additive}

Zscheischler et al.\ \citep{Zscheischler2020} demonstrate in the context of climate systems that compound risks --- where multiple hazard dimensions co-occur --- produce impacts that exceed the sum of individual-dimension risks. The mechanism in transportation networks is straightforward: the value of redundancy depends on the alternate being usable, and the value of hardening depends on the primary corridor being the only path.

Consider the counterfactuals:

\textbf{High B1, low D1}: I-90 Homestake Pass (Montana Continental Divide). Structurally isolated (B1 = 7.8); rarely closes due to winter precipitation (D1 = 6.4). This corridor is a contingency risk: if it closes, freight has no good alternate. But it closes only 25 times per year, with moderate duration. Investment in this corridor is justified but not urgent.

\textbf{Low B1, high D1}: I-95 South Carolina coast. Frequently exposed to hurricane surge and flooding (D1 = 6.8); well-connected with adequate alternates on US-17 and US-1 (B1 = 3.2). When I-95 South Carolina closes, freight reroutes with modest penalty. Infrastructure hardening investment here has local benefit but limited national-freight impact.

\textbf{High B1, high D1}: I-80 Donner Pass. Structurally isolated (B1 = 8.3) and frequently closes (D1 = 7.8). When Donner closes, there is no adequate alternate. The 50 annual closures are not independently manageable events --- each one is a national-freight disruption. The compound condition makes each D1 event a B1 event simultaneously.

The operational consequence is that compound exposure corridors produce disproportionate freight disruption cost per closure event. The data sources for this analysis are: ROUTE v1.2 corpus (227 corridors) \citep{ROUTE_A2}, Caltrans operational closure logs \citep{Caltrans2023}, FHWA incident database \citep{FHWA_freight2023}, and FEMA NFHL flood data \citep{FEMA_NFHL}.
